% ============================================
% 讲座开场
% ============================================
\section{讲座开场}

\subsection{自我介绍与讲座目标}

各位同学好，我是你们的朋辈学业导师。今天的讲座内容是 \textbf{C语言期末模拟试卷的讲评}。

本次讲座的目标：
\begin{itemize}
    \item 逐题讲解模拟试卷，帮助大家 \textbf{查漏补缺}
    \item 重点分析 \textbf{易错知识点} 和 \textbf{常见陷阱}
    \item 提炼 \textbf{解题思路与方法}，而非仅仅对答案
\end{itemize}

\textbf{教学提示：} 此处可简短介绍自己，拉近与学生的距离。强调本次讲解以「模拟卷」为线索，但目标是覆盖期末考试的核心考点。

\subsection{模拟考试分值分布}

本次模拟试卷满分 100 分，题型分布如下：

\begin{table}[h]
    \centering
    \begin{tabular}{@{}lll@{}}
        \toprule
        \textbf{题型} & \textbf{题量} & \textbf{分值} \\
        \midrule
        选择题   & 20题（每题2分） & 40分 \\
        基本概念填空题 & 10空（每空2分） & 20分 \\
        阅读程序填空题 & 10空（每空2分） & 20分 \\
        程序完善题 & 5空（每空2分） & 10分 \\
        编程题   & 1题          & 10分 \\
        \bottomrule
    \end{tabular}
\end{table}

从分值来看，选择题占 40 分，是拿分的大头；填空题和阅读题各占 20 分，考验基本功；编程题 10 分，区分度最高。

\subsection{学生常见薄弱环节}

根据往年考试和本次模拟的情况，以下几个知识点学生 \textbf{丢分最多}：

\begin{enumerate}
    \item \textbf{指针} —— 指针与数组的关系、指针作为函数参数、函数指针声明。很多同学「听得懂但做不对」。
    \item \textbf{递归} —— 递归调用过程的理解、递归与循环的转换。建议画调用栈图辅助理解。
    \item \textbf{宏定义展开} —— 宏是纯文本替换，不是函数调用！特别注意括号问题（如 \texttt{\#define SQR(x) x*x} 的陷阱）。
    \item \textbf{二维数组} —— 初始化方式、存储顺序（行优先）、下标范围。
    \item \textbf{逗号表达式} —— 从左到右依次求值，取最右值为结果。
\end{enumerate}

\textbf{教学提示：} 此处可追问学生：「你们觉得哪些题最难？」根据学生反馈适当调整后续讲解的详略。

\subsection{讲座安排}

本次讲评按试卷顺序进行，分为以下几个板块：

\begin{enumerate}
    \item 选择题 Q1--Q10（标识符、优先级、字面量、ASCII、逻辑运算、break/continue、printf、数组声明与越界）
    \item 选择题 Q11--Q20（函数、宏、指针、联合体、结构体、循环、static、字符串）
    \item 基本概念填空题（源文件后缀、头文件、三种结构、取地址、复合赋值、函数声明、二维数组初始化、逗号表达式）
    \item 阅读程序填空题（循环控制、指针参数、数组反转、字符处理、递归）
    \item 程序完善题（矩阵运算、素数判断）
    \item 编程题（电梯调度系统）
\end{enumerate}

讲解过程中，对每道题我会先点明 \textbf{考查的知识点}，再讲 \textbf{解题思路}，最后提醒 \textbf{常见错误}。欢迎大家随时打断提问。

\subsection{参考建议}

\begin{itemize}
    \item 听讲时重点记 \textbf{「我当时为什么做错了」}，而不是抄答案
    \item 讲座结束后，建议把错题 \textbf{重新做一遍}，确保真正理解
    \item 编程题一定要 \textbf{上机练习}，光看不练是学不会的
\end{itemize}

好，下面我们开始讲评。先从选择题 Q1--Q10 开始。
